\section{Quintic B-spline collocation}\label{sec:method}

\subsection{Spline spaces and grids}

We approximate $w(t,\cdot,\cdot)$ in a tensor-product space of quintic splines,
\begin{equation}\label{eq:tensor}
w(t,x,v) \approx \sum_{i=1}^{n_x}\sum_{j=1}^{n_v} c_{ij}(t)\,N_i(x)\,M_j(v),
\end{equation}
where $N_i$ and $M_j$ are B-splines of degree $p=5$ on open knot vectors (end knots repeated $p+1$ times)
\citep{deBoor2001}. The knots are the nodes of stretched grids. In each direction the nodes are
$y^\ast + \alpha\sinh(s)$ with $s$ uniform, so they cluster around a point $y^\ast$ with a width set by
$\alpha$. In $x$ we cluster around $x^\ast=0$ with $\alpha_x=0.3$: this is where the payoff has its kink
and where the $x$-diffusion vanishes. In $v$ we cluster around $v^\ast=0$ with $\alpha_v=0.15$.
For the convergence study of Section~\ref{sec:convergence} we fix the map exactly: the $s$-interval is
split at $s=0$ (the image of $x=0$) and the two sides receive $6k$ and $11k$ uniform steps, so that
$N_x = 17k+1$ and every grid in the sequence is an image of the same piecewise-linear map in $s$. Without
this precaution the number of nodes on each side of $x=0$ jumps by rounding from one grid to the next,
and the error constant jumps with it.

The payoff \eqref{eq:terminal} is only continuous at $x=0$, while a quintic spline with simple knots is
$C^4$. We insert $x=0$ as a knot of multiplicity four (removing any node closer than a quarter of the
local spacing), which makes the space $C^1$ there. A multiplicity of five would allow a kink in the
spline itself, but the second derivative would then be undefined at the kink, and a collocation method for
a second-order equation needs it. With multiplicity four the interpolant of the payoff is not exact, but
the error is confined to a few intervals around the kink.

The collocation points are the Greville abscissae $\xi_k = (t_{k+1}+\dots+t_{k+p})/p$ of each knot
vector. There are as many of them as there are basis functions, and they include both end points. In
particular $v=0$ is a collocation line.

\subsection{Collocation equations}

Let $B^{(d)}_x$ be the $n_x\times n_x$ matrix with entries $N_i^{(d)}(\xi_k)$, $d=0,1,2$, and define
$B^{(d)}_v$ in the same way. On the tensor grid of collocation points, and with $c$ the vector of
coefficients ordered by $x$ first, the values and derivatives of \eqref{eq:tensor} are
\begin{equation*}
M = B^{(0)}_x\otimes B^{(0)}_v,\quad
D_{xx} = B^{(2)}_x\otimes B^{(0)}_v,\quad
D_{xv} = B^{(1)}_x\otimes B^{(1)}_v,\quad
D_{vv} = B^{(0)}_x\otimes B^{(2)}_v,
\end{equation*}
and similarly $D_x$, $D_v$. Each one-dimensional matrix has at most six nonzeros per row, so the
Kronecker products have at most 36. Writing $\operatorname{diag}(f)$ for the diagonal matrix of a
coefficient $f$ evaluated at the collocation points, the spatial operator of \eqref{eq:pde} becomes
\begin{equation}\label{eq:Lh}
L = \operatorname{diag}\bigl(\tfrac12 v x^2\bigr)D_{xx} - \operatorname{diag}(\rho\xi v x)D_{xv}
+ \operatorname{diag}\bigl(\tfrac12\xi^2 v\bigr)D_{vv} - \operatorname{diag}\bigl(rx+\tfrac1T\bigr)D_x
+ \operatorname{diag}\bigl(\kappa(\theta-v)+\rho\xi v\bigr)D_v .
\end{equation}
The mixed derivative is simply one more sparse term; no splitting is involved.

Boundary conditions replace collocation equations. On the lines $x=x_{\min}$ and $x=x_{\max}$ and on
$v=v_{\max}$ the rows of \eqref{eq:Lh} are replaced by rows that evaluate \eqref{eq:bcx} or \eqref{eq:bcv}
with the exact derivatives of the basis functions. On $v=0$ nothing is replaced. There the diagonal
factors of the three second-order terms in \eqref{eq:Lh} are zero, and the remaining row is the
collocation of the degenerate equation \eqref{eq:v0}. The Fichera analysis of
Section~\ref{sec:boundaries} is thus built into the scheme without any special treatment, and the same
holds for the geometric twin.

\subsection{Time stepping}

In time-to-maturity $\tau=T-t$ the semi-discrete problem is $M\dot c = Lc$ (with boundary rows held at
their homogeneous values). We use the $\theta$-scheme
\begin{equation}\label{eq:theta}
(M - \theta\Delta\tau\,L)\,c^{n+1} = (M + (1-\theta)\Delta\tau\,L)\,c^n ,
\end{equation}
with four backward Euler steps ($\theta=1$) at the start and Crank--Nicolson ($\theta=\tfrac12$)
afterwards. The implicit start-up \citep{Rannacher1984} damps the high-frequency error created by the
kink; without it Crank--Nicolson is known to lose accuracy on non-smooth data
\citep{PooleyVetzalForsyth2003}. The initial coefficients $c^0$ are obtained by collocating the payoff,
$Mc^0 = w(T,\cdot,\cdot)$ at the collocation points, with the boundary rows included.

Since $L$ does not depend on $\tau$, the two matrices on the left of \eqref{eq:theta} are factorised once
each by a sparse LU decomposition, and each time step costs a sparse matrix--vector product and two
triangular solves. On the grid used for most results ($N_x=273$ nodes, $N_v=61$) the system has
$n_xn_v = 280\times65 = \num{18200}$ unknowns.

The error of Crank--Nicolson with Rannacher start-up behaves like
$c_2\Delta\tau^2 + c_3\Delta\tau^3 + \dots$; the cubic term comes from the four implicit steps, whose
local error is not symmetric in time. Where we need the time error to be negligible we therefore use two
levels of Richardson extrapolation,
\begin{equation}\label{eq:richardson}
R_1(\Delta\tau) = \frac{4P(\Delta\tau/2) - P(\Delta\tau)}{3},\qquad
R_2(\Delta\tau) = \frac{8R_1(\Delta\tau/2) - R_1(\Delta\tau)}{7},
\end{equation}
which needs three solves with $N_t$, $2N_t$ and $4N_t$ steps. Section~\ref{sec:convergence} confirms both
the second-order behaviour and the cubic remainder.

\subsection{Reading the price}

The price follows from \eqref{eq:readout} by evaluating the spline \eqref{eq:tensor} at
$(K/S_0, v_0)$:
\begin{equation*}
C(K) = S_0\sum_{i,j} c_{ij}\,N_i(K/S_0)\,M_j(v_0).
\end{equation*}
No interpolation between grid values is involved, and the evaluation is repeated for as many strikes as
needed at negligible cost. This removes a source of error that is easy to overlook with finite
differences, where the price at an off-grid point has to be interpolated and the interpolation error adds
to the discretisation error.

\subsection{The geometric twin}

The same construction applies to \eqref{eq:pdegeo}, with the kink knot at $\zeta=\ln K$. Its operator
depends on $\tau$ only through $a=\tau/T$ and splits as
\begin{equation*}
L(\tau) = a^2 M_2 + a M_1 + M_0 ,
\end{equation*}
where $M_2$ collects the $\zeta\zeta$ term, $M_1$ the $\zeta$ and $\zeta v$ terms and $M_0$ the rest. The
three matrices are assembled once, but the system matrix changes at every step and must be factorised
again. This is the main cost difference between the two problems and it matters for the comparison with
ADI in Section~\ref{sec:adi}.
